\section{Background and Motivation}
\label{sec:problem}

\subsection{Periodic Execution and Retained State}
\label{sec:request-classes}
\label{sec:formulation}

A periodic interaction session carries its history from one micro-turn to the next.
The application determines which history to retain, whether through full retention, a fixed window, or a summary.
Residency management changes where the corresponding KV is stored while preserving that choice.
The opportunity to release GPU space depends on when the model uses the state: periodic input capture alone does not establish periodic model execution or an interval without KV accesses.

For timing analysis, consider sessions with a common period $T$ and reference time $t_0$.
Session $i$ has phase $\phi_i\in[0,T)$, so the tick that releases micro-turn $k$ is
\begin{equation}
  r(i,k)=t_0+\phi_i+kT.
  \label{eq:release}
\end{equation}
Each micro-turn should finish generation by the next tick:
\begin{equation}
  c(i,k)\le r(i,k+1)=r(i,k)+T,
  \label{eq:deadline}
\end{equation}
where $c(i,k)$ is generation completion.
This is a soft deadline, with occasional misses allowed under a specified miss-rate limit.
Latency is measured from the planned tick as $c(i,k)-r(i,k)$; delaying input submission or computation does not move the deadline.
Delivery latency and initial phase-alignment delay require separate accounting (Section~\ref{sec:methodology}).
\placeholder{Allowed deadline-miss rate, observation interval, and model-specific generation-completion event.}

Let $q(i,k)$ be the last use of the KV selected for eviction in micro-turn $k$.
That state can leave the GPU once outstanding references are released and a valid host copy is available.
It must return before the next computation accesses it; \systemname targets completion by $r(i,k+1)$ to keep restoration off the next micro-turn's execution path.
The interval from $q(i,k)$ to that tick is the planned window for reclamation and restoration, provided there are no intervening accesses.
Late computation can shorten or eliminate this window.

\subsection{When KV Capacity Limits Concurrency}
\label{sec:capacity-problem}
\label{sec:slack-conditions}

A session's KV-idle interval does not imply spare compute capacity: other sessions may keep the GPU busy throughout it.
Nor do periodic releases alone establish compute slack~\cite{liu1973scheduling}.
The relevant question is whether \emph{all admitted sessions} still leave time within the period when their retained KV fills GPU memory.
Even if each micro-turn adds a predictable amount of state, attention to a growing history can increase execution time.
The memory and compute limits must therefore be compared at the same context length.

Consider a reference configuration with $N$ sessions whose inputs are available at the same tick and whose period is $T$.
Each session reaches a maximum retained length of $L$ positions within the period, with $\kappa$ KV bytes per position and no inter-session sharing.
Let $G_{\mathrm{KV}}$ be the GPU capacity available for KV after weights, activations, and other runtime allocations.
Ignoring allocation rounding, full residency fills this pool at
\begin{equation}
  L_{\mathrm{mem}}=\frac{G_{\mathrm{KV}}}{N\kappa}.
  \label{eq:memory-boundary}
\end{equation}
Let $C_N(L)$ be the elapsed time needed for all $N$ sessions' per-period work under a specified execution policy with the required KV resident.
It includes input processing, model execution, and scheduling overhead; batching means that it is not generally the sum of isolated session times.

The time left at the memory limit is $T-C_N(L_{\mathrm{mem}})$.
To bound it, hold the per-period input and generation work fixed and suppose that, over the context range of interest,
\begin{equation}
  C_N(L)\le A_N+B_NL,\qquad A_N\ge0,\ B_N\ge0.
  \label{eq:compute-envelope}
\end{equation}
The intercept $A_N$ accounts for work independent of history length; the slope $B_N$ allows execution cost to grow with that length.
This affine upper bound is an assumption to establish for the chosen workload, platform, and execution policy.
If $L_{\mathrm{mem}}$ lies in the same range, the remaining time satisfies
\begin{equation}
  T-C_N(L_{\mathrm{mem}})
  \ge T-A_N-\frac{B_NG_{\mathrm{KV}}}{N\kappa}.
  \label{eq:memory-before-compute}
\end{equation}
A positive right-hand side is sufficient for memory to fill before the period's compute budget is exhausted.
If the right-hand side is at least $\eta T$ for $0<\eta<1$, all $N$ sessions leave at least a fraction $\eta$ of the period unused in this reference configuration.

Coefficients fitted to measurements yield a prediction of slack, not a proven lower bound, and require validation on independent configurations.
Even an established slack bound does not show that an additional session or recomputation will meet its deadline: the extra work, recovery costs, and interference must fit while preserving each session's deadline.
\placeholder{Motivating measurement: joint KV allocation and aggregate execution time at the full-residency limit, with concurrency and context length; validate the predicted compute slack on separate configurations.}

\subsection{Restoration Needs Both Time and Space}
\label{sec:opportunity}

Removing a session's KV frees capacity only until the server allocates space to restore it.
The full restoration destination must be allocated before copying starts, even though its KV is not yet usable.
Let $\mathcal{P}(t)$ be the set of physical GPU blocks allocated to KV at time $t$, including destinations of unfinished restorations, and let $b_p$ be the size of block $p$.
A feasible schedule requires
\begin{equation}
  \sum_{p\in\mathcal{P}(t)} b_p \le G_{\mathrm{KV}}
  \quad\text{at every }t,
  \label{eq:memory}
\end{equation}
with each shared physical block counted once.
Counting only valid KV would miss the allocation occupied by ongoing restorations.

That allocation must also occur early enough for the copy to finish.
For a restoration of $e$ bytes starting at time $u$, let $\widehat B$ be its estimated effective host-to-device bandwidth and $\delta$ an allowance for scheduling and completion overhead.
The estimated completion must satisfy
\begin{equation}
  u + e/\widehat B + \delta \le r(i,k+1).
  \label{eq:restore-time}
\end{equation}
The estimate must reflect bandwidth available under the shared transfer schedule, including the effects of host backup traffic and concurrent GPU work, rather than the link's nominal rate.

These constraints pull restoration timing in opposite directions.
An earlier restoration has more time to finish but holds its destination longer; a larger eviction releases more memory but takes longer to restore.
Session phases change when space becomes available and when restored state is needed.
A schedule can consequently fit its total transfer bytes within a period and still fail because transfers overlap or their destinations exceed GPU capacity at some instant.
Memory multiplexing requires checking both the transfer windows and GPU allocation over time, together with computation under the chosen phases.
